\documentclass[12pt]{article}
\usepackage[utf8]{inputenc}
\usepackage[T1]{fontenc}
\usepackage{hyperref}
\title{Tringgus-Sembaan Bidayuh Language Resources for K-12 Curriculum Development}
\author{Compiled by Claude Sonnet 5 (claude-sonnet-5) for Andr\'as Kornai}
\date{}
\begin{document}
\maketitle
\section{Introduction}
Tringgus-Sembaan Bidayuh (ISO 639-3 trx) is a Land Dayak (Austronesian) language of the Tringgus area of Padawan, Sarawak, Malaysia. Almost nothing specific to this language was found in the searches.
\section{Primary Educational Resources}
None were found.
\section{Linguistic and Reference Materials}
\begin{itemize}
\item Glottolog entry for Tringgus: \url{https://glottolog.org/resource/languoid/id/trin1272.bigmap.html}
\item Glottolog reference: Rensch, Rensch, Noeb and Ridu 2012: \url{https://glottolog.org/resource/reference/id/318599}
\item Overview of the Land Dayak languages: \url{https://en.wikipedia.org/wiki/Land_Dayak_languages}
\end{itemize}
\section{Religious and Community Resources}
Joshua Project lists Bible translation as ``Translation Started'', with no portions published: \url{https://joshuaproject.net/index.php/languages/trx} and \url{https://joshuaproject.net/people_groups/15577/MY}
\section{Digital and Audio Resources}
None were found. Joshua Project points to directories such as Global Recordings Network and SIL resources, which were not checked.
\section{Curriculum Development Context}
Joshua Project estimates about 1,600 speakers. A 2026 study of Bidayuh in Sarawak generally (not Tringgus specifically) reports that most children no longer speak the language. Any curriculum would have to start from field documentation.
\end{document}